% Paper 6: C-H Activation — Figure Captions

\begin{figure}[htbp]
\caption{\textbf{Substrate Positioning at the Eve of C--H Activation (Paper 6, Panel 1).}
(\textbf{A})~Activation partition depths $\Delta\mathcal{M}_{\mathrm{HAT}}$ for five
substrate classes (aliphatic, benzylic, allylic, aromatic, epoxidation). Aliphatic
C--H has the highest barrier ($\Delta\mathcal{M} = 0.65$); epoxidation the lowest
($0.35$).
(\textbf{B})~S-coordinates $S_k$, $S_t$, $S_e$ for Compound~I across the five
transition states; $S_k$ tracks the charge-character change associated with increasing
activation.
(\textbf{C})~S-distance $d_S(\mathrm{Cpd\,I},\,\mathrm{TS})$ per substrate class.
Distance and $\Delta\mathcal{M}$ are proportional, confirming the S-entropy coordinate
as a reliable activation-depth proxy.
(\textbf{D})~Three-dimensional landscape of $E_a$ versus substrate class and S-distance,
illustrating the full parameter space sampled by P450 C--H activation reactions.}
\label{fig:substrate_positioning}
\end{figure}

\begin{figure}[htbp]
\caption{\textbf{The C--H Bond-Order Coordinate and HAT Three-Body Trajectory (Panel 2).}
(\textbf{A})~Binary C--H bond-order states: $\beta_{\mathrm{C-H}} = 1$ (intact, green)
and $\beta_{\mathrm{C-H}} = 0$ (cleaved, red). The binary character mirrors the O--O
coordinate introduced in Paper~5.
(\textbf{B})~TS geometry: C--H distance (1.30~\AA{}), O--H distance (1.20~\AA{}),
and C--H--O linearity angle (165°) from QM/MM studies of CYP3A4
\citep{shaik2010}, re-expressed as a single partition-cell boundary.
(\textbf{C})~Decomposition of $\Delta\mathcal{M}_{\mathrm{HAT}} = 0.65$ into two
contributions: half-C--H cleavage (0.347) and Fe partial reduction at the TS (0.30).
(\textbf{D})~Three-dimensional $\Delta\mathcal{M}_{\mathrm{HAT}}$ landscape over
substrate BDE (80--120~kcal/mol) and C--H--O approach angle (140--180°). The CYP3A4
aliphatic operating point is marked in red.}
\label{fig:hat_coordinate}
\end{figure}

\begin{figure}[htbp]
\caption{\textbf{Kinetic Isotope Effect: ZPE and Tunneling Contributions (Panel 3).}
(\textbf{A})~Decomposition of the predicted KIE $\approx 7.2$: classical ZPE
contribution ($\approx 6.2$, orange) and tunneling factor $\kappa_H/\kappa_D \approx
1.16$ (green), within the experimental range 4--11 (grey band).
(\textbf{B})~$\Delta$ZPE (kcal/mol) versus C--H stretch frequency~(cm$^{-1}$); the
shaded band indicates the typical P450 substrate range (2800--3100~cm$^{-1}$).
(\textbf{C})~KIE versus temperature (250--400~K). The ZPE component decreases with
$T$; the tunneling factor is nearly temperature-independent. At 310~K, KIE~$\approx
7.2$.
(\textbf{D})~Three-dimensional KIE surface over temperature and $\Delta$ZPE,
with the CYP3A4 operating point (red marker).}
\label{fig:kie}
\end{figure}

\begin{figure}[htbp]
\caption{\textbf{The Substrate-Radical Intermediate: Partition Cell and Lifetime (Panel 4).}
(\textbf{A})~Partition-cell occupancy (0/1) for $\beta_{\mathrm{C-H}}$,
$\beta_{\mathrm{O-H}}$, $\beta_{\mathrm{C-O}}$, and $\sigma_{\mathrm{rad}}$ in the
reactant (Cpd~I + substrate), radical intermediate, and product states.
(\textbf{B})~Rebound rate $\krebound = 7.4 \times 10^9$~s$^{-1}$ (red line) vs
escape rates $k_{\mathrm{escape}} = 10^8$--$5 \times 10^9$~s$^{-1}$ (blue points);
log scale.
(\textbf{C})~Stereoretention $f_{\mathrm{retained}}$ versus $k_{\mathrm{escape}}$ on
a log scale. The shaded band (40--90\%) marks the experimental range; the framework
prediction crosses the band for $k_{\mathrm{escape}} = 10^8$--$10^9$~s$^{-1}$.
(\textbf{D})~Three-dimensional $f_{\mathrm{retained}}$ landscape over
$(\log_{10}(k_{\mathrm{rebound}}),\, \log_{10}(k_{\mathrm{escape}}))$.}
\label{fig:radical_intermediate}
\end{figure}

\begin{figure}[htbp]
\caption{\textbf{Oxygen-Rebound Aperture: Bond Formation and Stereospecificity (Panel 5).}
(\textbf{A})~Activation depth comparison: $\Delta\mathcal{M}_{\mathrm{HAT}} = 0.65$
(blue), $\Delta\mathcal{M}_{\mathrm{rebound}} = 0.30$ (green), difference
$\Delta\Delta\mathcal{M} = 0.35$ (orange). Rebound is intrinsically faster by a factor
$e^{0.35} \approx 1.4$.
(\textbf{B})~Rate comparison with literature: $\kHAT$ (predicted), $\krebound$
(predicted), Newcomb 1995 lower bound ($> 10^9$~s$^{-1}$), and Groves 1985 estimate
($\sim 10^{10}$~s$^{-1}$).
(\textbf{C})~C--O bond formation trajectory: $\beta_{\mathrm{CO}}$ transitions from 0
(no bond) to 1 (C--OH product) through the rebound TS at $\tau = 1.5 \tau_p$.
(\textbf{D})~Three-dimensional $\krebound$ landscape over
$(\Delta\mathcal{M}_{\mathrm{rebound}},\, \xi_{\mathrm{rad}})$, with the operating
point at $(0.30,\, 0.57)$ marked in red.}
\label{fig:rebound}
\end{figure}

\begin{figure}[htbp]
\caption{\textbf{Testosterone Regioselectivity: CYP3A4 6$\beta$-Hydroxylation (Panel 6).}
(\textbf{A})~Effective rates $k_i^{\mathrm{eff}} = \nu_{\mathrm{floor}} \cdot g_i \cdot
e^{-\Delta\mathcal{M}_i}$ for four competing testosterone positions
($6\beta,\, 2\beta,\, 15\beta,\, 16\beta$). The $6\beta$ position dominates.
(\textbf{B})~Predicted regioselectivity fractions (pie chart): $6\beta\ 49\%$,
$16\beta\ 23\%$, $2\beta\ 17\%$, $15\beta\ 12\%$. Experimental dominant product
is $6\beta$ at 50--70\% in human liver microsomes \citep{guengerich1998}.
(\textbf{C})~Sensitivity of $f_{6\beta}$ to the geometric factor $g_{6\beta}$
(sweep 0.5--2.0). The prediction exceeds 45\% for all $g_{6\beta} > 0.9$.
(\textbf{D})~Three-dimensional selectivity landscape over
$(\Delta\mathcal{M}_{6\beta},\, g_{6\beta})$.}
\label{fig:regioselectivity}
\end{figure}

\begin{figure}[htbp]
\caption{\textbf{Five Reaction Types Unified Under the Three-Body Aperture (Panel 7).}
(\textbf{A})~Activation depths $\Delta\mathcal{M}_{\mathrm{HAT}}$ per class:
aliphatic (0.65) $>$ benzylic (0.50) $>$ allylic (0.45) $>$ aromatic (0.38)
$>$ epoxidation (0.35). Monotonic trend reflects increasing substrate activation.
(\textbf{B})~Intrinsic rates $k = 10^{10} \cdot e^{-\Delta\mathcal{M}}$ per class on
a log scale.
(\textbf{C})~KIE predictions: aliphatic $\approx 7.2$, benzylic $\approx 5.0$,
allylic $\approx 4.4$; aromatic and epoxidation KIE $= 1$ (no primary H isotope
effect, since the reaction coordinate does not involve H motion).
(\textbf{D})~Three-dimensional landscape of reaction class index, activation depth,
and intrinsic rate.}
\label{fig:reaction_types}
\end{figure}

\begin{figure}[htbp]
\caption{\textbf{Full Validation Summary: Eight Independent Observables (Panel 8).}
(\textbf{A})~Predicted versus experimental scatter for eight observables
($\Delta\mathcal{M}_{\mathrm{HAT}}$, $E_a$, $\kHAT$, $\krebound$,
$\krebound/\kHAT$, KIE, stereoretention, $6\beta$-regioselectivity). Diagonal $y=x$
reference line; green points fall within 20\% of experimental values.
(\textbf{B})~Per-observable relative error (horizontal bars). All observables are
within 30\% of their experimental reference values.
(\textbf{C})~Pie chart: 8/8 observables PASS at the 30\% threshold.
(\textbf{D})~Three-dimensional bar landscape: predicted (orange) versus normalised
experimental reference (blue) for all eight observables.}
\label{fig:validation}
\end{figure}
